% Generated draft: requires mathematical and source review.
\documentclass[11pt]{article}
\usepackage[margin=25mm]{geometry}
\usepackage{fontspec}
\setmainfont{Latin Modern Roman}
\usepackage{amsmath,amssymb,mathtools,unicode-math}
\setmathfont{Latin Modern Math}
\usepackage{enumitem,needspace,xurl,hyperref}
\setlength{\parindent}{0pt}
\setlength{\parskip}{5pt}
\setlength{\emergencystretch}{3em}
\newcommand{\N}{\mathbb{N}}
\newcommand{\Z}{\mathbb{Z}}
\newcommand{\Q}{\mathbb{Q}}
\newcommand{\R}{\mathbb{R}}
\newcommand{\C}{\mathbb{C}}
\newcommand{\F}{\mathbb{F}}
\newcommand{\A}{\mathbb{A}}
\newcommand{\PP}{\mathbb{P}}
\newcommand{\E}{\mathbb{E}}
\newcommand{\eps}{\varepsilon}
\newcommand{\dd}{\,\mathrm{d}}
\begin{document}
\section*{UnsolvedMath ID: 30001709}
\textbf{Source problem code:} catalogue entry
\subsection{Definitions and mathematical statement}
Write
\[
M:=M_G(\Sigma;C_1,\ldots,C_r)
\]
for the quasi-Hamiltonian \(G\)-space named in the supplied record. A quasi-Hamiltonian structure involves a \(G\)-action, an invariant two-form, and a group-valued moment map \(\Phi:M\to G\), subject to the quasi-Hamiltonian axioms. Here \emph{pre-quantized} means that a compatible prequantization has been supplied; it does not assert that every such space admits one. Write \(Q(M)\) for the quantization requested by the source.

\textbf{Problem as recorded.} Given a pre-quantized quasi-Hamiltonian \(G\)-space \(M_G(\Sigma;C_1,\ldots,C_r)\), compute
\[
Q\bigl(M_G(\Sigma;C_1,\ldots,C_r)\bigr).
\]
The request particularly concerns non-simply connected \(G\), for which the source notes obstructions to prequantization. The computation is requested only for data equipped with a prequantization; the record does not assert that those obstructions always vanish.

\textbf{Unresolved qualifications.} The supplied record does not specify the admissible class of groups \(G\), the topology, orientation, compactness, or boundary conditions of \(\Sigma\), the domain of \(r\), or the meaning and admissibility conditions of the parameters \(C_i\). In particular, interpreting the \(C_i\) as conjugacy classes prescribing boundary holonomy would require information absent from this record. It also does not define the construction of \(M_G(\Sigma;C_1,\ldots,C_r)\), the normalization and level of its prequantization, or the definition and target of \(Q\). Thus the displayed request preserves the source formulation, but a fully specified mathematical computation requires these additional data. No particular moduli-space construction, quantization convention, or range of boundary cases is selected here.
\subsection{Short English statement}
Compute the quantization of the source-designated pre-quantized quasi-Hamiltonian moduli space, particularly when \(G\) is not simply connected. The group, surface, boundary data, prequantization conventions, and meaning of quantization require clarification before this becomes a fully specified problem.
\subsection{Sources}
\begin{itemize}
\item Local UnsolvedMath registry, permanent ID \(30001709\), source problem code \texttt{catalogue entry}, title \emph{Quantization of Quasi-Hamiltonian Moduli Spaces}; collection: \emph{Oberwolfach Reports Open Problems}. The supplied record gives proposed year 2011 and does not identify an individual proposer.
\item \emph{Geometric Quantization in the Non-compact Setting} (2011), Oberwolfach Reports, pp.\ 425--521, DOI 10.4171/owr/2011/09. Supplied report citation and original-report literature source: \url{https://doi.org/10.4171/owr/2011/09}. The record describes Oberwolfach Reports as documenting mathematical workshops and their open-problem discussions; it gives no precise location of this problem within the report.
\item Meinrenken, \emph{Verlinde formulas for nonsimply connected groups}. Supplied literature source: \url{https://arxiv.org/abs/1706.04045}.
\item Supplied literature assessment, checked on \texttt{2026-08-21}: the record labels the problem \texttt{partially\_solved} and states that quantization is established for simply connected groups, while later work proves important nonsimply connected Verlinde cases, including surfaces with at most one boundary component; it describes the full problem as broader. This is supplied dated curation derived from the OWR collection, not an independently verified literature assessment. The record explicitly calls for independent verification of time-sensitive status and cited claims.
\end{itemize}
\end{document}
